\documentclass[11pt]{article}
\emergencystretch=2em
\usepackage[margin=1in]{geometry}
\usepackage[T1]{fontenc}
\usepackage{lmodern,amsmath,amssymb,amsthm,booktabs,array,enumitem}
\usepackage[colorlinks=true,linkcolor=blue,citecolor=blue,urlcolor=blue]{hyperref}
\newtheorem{theorem}{Theorem}
\newtheorem{proposition}[theorem]{Proposition}
\newtheorem{corollary}[theorem]{Corollary}
\newtheorem{lemma}[theorem]{Lemma}
\theoremstyle{definition}
\newtheorem{definition}[theorem]{Definition}
\newtheorem{example}[theorem]{Example}
\theoremstyle{remark}
\newtheorem{remark}[theorem]{Remark}
\newcommand{\Prb}{\mathbb P}
\newcommand{\E}{\mathbb E}
\newcommand{\ind}{\mathbf 1}
\newcommand{\pr}{\mathrm{pr}}
\newcommand{\conv}{\mathrm{conv}}
\title{Allocation and Coordination under Scarce Resources:\\Admissible Routing, Prices, Review Sets,\\and Containment for Gated Work}
\author{Ivo Matija\v{s}evi\'c}
\date{October 2026}
\newcommand\blfootnote[1]{\begingroup\renewcommand\thefootnote{}\footnote{#1}\addtocounter{footnote}{-1}\endgroup}
\begin{document}
\maketitle
\blfootnote{\textcopyright~2026 Ivo Matija\v{s}evi\'c.\ This work is licensed under a Creative Commons
Attribution 4.0 International (CC BY 4.0) license, \url{https://creativecommons.org/licenses/by/4.0/}.\ Preprint, DOI \href{https://doi.org/10.5281/zenodo.23449199}{10.5281/zenodo.23449199}.}
\begin{abstract}
Earlier papers in this series decide which executor, gate stack, or mechanism is admissible for one
step of work. They do not decide how scarce capacity is shared when many classes of work compete for
it: a principal's attention, verification memory, external review lineages, rate-limited seats, and
calendar time. This paper studies that allocation problem. We allocate only among options whose
safety premises remain valid under the allocation actually chosen. We keep three distinctions: fluid
feasibility versus executable schedules, marginal prices versus finite policy changes, and measured
review value versus worst-case review value. In a fluid linear model, the feasible and optimal
allocations do not depend on the executor's error law when bounds decide admissibility and the
resource use, costs, demands, and capacities are law-independent or replaced by common uniform bounds.
Feasibility has an exact certificate by nonnegative resource prices; subset cuts are exact for
single-resource pools but insufficient for general resource bundles. Shadow prices value the principal's minutes, but a new option strictly lowers the optimum
exactly when no old optimal dual prices it out, so a negative reduced cost at one dual is not enough.
Choosing optional principal reviews is a knapsack. Review value under a fixed fault law is submodular;
worst-case review value need not be, and a reviewer can have zero guaranteed value alone and positive
guaranteed value in a coalition. Verifying an implementation does not identify the intended
specification. Interface contracts bound severity: parallel escape routes add, and serial stages
multiply only through conditional ceilings. With optional participation, achievable throughput does not
fall as workers are added. Substituting a mechanization for scarce work needs feasibility in every
resource and every period. Queueing models, a scheduling formulation of verification compute, and
MaxWeight stability complete the picture, each under its stated assumptions. Many results are
classical; we state the conditions under which they apply to gated work. These are mathematical
models, not measurements.
\end{abstract}

\section{Introduction}
Part~I of this series shows that a gate family with a common blind set leaves every executor with a
verification floor~\cite{p1}. Part~II shows that gates generated within one family inherit its blind
set~\cite{p2}. Part~III treats escape, cost, and correlated escape~\cite{p3}. Part~IV studies the move
of work from executor calls to deterministic mechanisms~\cite{p4}. Part~V removes the
constant-fault-rate assumption~\cite{p5}. Part~VI adds
severity~\cite{p6}. It defines when an option is admissible for a step: its worst-case expected loss,
over the models compatible with the evidence, fits a severity-derived budget. That decision uses only
floors, ceilings, and certified levels, so it holds for every error law within those bounds.

These papers decide one step at a time. A working system runs many classes of work at once, and they
compete for capacity. Some resources are scarce: the principal's attention, the memory of a
verification machine, the capacity of an external review lineage, rate-limited seats, and calendar
time. This paper asks which worker or mechanism should receive which unit when such resources bind,
and how the organization of the work (decomposition, interfaces, parallelism) changes risk and
capacity.

The answer rests on one rule and three distinctions. The rule: allocate only among options whose
admissibility holds under the allocation actually chosen, including its concurrency and review
arrangement. The distinctions: fluid feasibility versus executable schedules; marginal prices versus
finite policy changes; and measured review value versus worst-case review value.

\paragraph{Contributions.}
\begin{enumerate}[leftmargin=*,itemsep=2pt]
\item A fluid allocation model over admissible options, with the qualifications that admissibility
needs (Definition~\ref{def:model}), and its invariance to the executor's error law
(Theorem~\ref{thm:invariance}).
\item Capacity and feasibility: the bottleneck bound, Hall's condition for single-resource pools, and
an exact price certificate for resource bundles (Proposition~\ref{prop:bottleneck},
Theorems~\ref{thm:hall} and~\ref{thm:bundle}).
\item Shadow prices with the exact condition under which a new option helps, and a counterexample to
the unrestricted reduced-cost test (Theorem~\ref{thm:prices}, Example~\ref{ex:degenerate}); few split
classes at a fluid optimum (Proposition~\ref{prop:split}).
\item The principal's attention: a knapsack with the conditions it needs
(Proposition~\ref{prop:knapsack}); submodular fixed-law review value and non-submodular robust review
value (Theorem~\ref{thm:submodular}); the principal as producer (Corollary~\ref{cor:blind}); and the
gap between verification and intent (Proposition~\ref{prop:intent}).
\item Containment: modularity and serial containment bounds on severity
(Theorems~\ref{thm:modularity} and~\ref{thm:serial}); parallelism on coupled work
(Proposition~\ref{prop:usl}); review lineages as capacity (Corollary~\ref{cor:lineages}).
\item Substitution and scheduling: per-resource, per-period substitution feasibility
(Proposition~\ref{prop:lifecycle}); queueing at a scarce gate (Proposition~\ref{prop:queue}); a
scheduling formulation of verification compute (Proposition~\ref{prop:scheduling}); and MaxWeight
stability over admissible configurations (Theorem~\ref{thm:maxweight}).
\end{enumerate}

\paragraph{Scope and status of the results.}
A result whose heading carries a citation is a known result. We state it in this paper's notation,
with the conditions that our use needs. Every other result is proven here in the stated model; that
label does not claim originality. No result in this paper is a conjecture. The contribution is the
set of conditions under which the classical tools apply to gated work, and the counterexamples that
mark where they fail.

\section{Model}\label{sec:model}
\paragraph{Background from the series.}
A unit of work is produced by an executor and passes a stack of gates. An \emph{escape} is a faulty
unit that passes every gate. An executor's conditional escape probability, given the full history,
has a \emph{floor} and a \emph{ceiling}. A process whose conditional escape lies between the bounds
$\ell$ and $u$ at every history belongs to the class $\mathcal P(\ell,u)$; mixtures of executors and
adaptive routing among them stay in the class~\cite{p6}. A gate accepts every fault in its blind set.
A blind set is \emph{complete} if it is exactly the set of accepted faults. A conjunctive review stack
accepts a fault only if every member accepts it, so with complete blind sets a producer with badness
$p_e$ and fault law $\mu_e$ escapes the reviewers $r_1,\dots,r_k$ with probability
$p_e\mu_e(\bigcap_j\Sigma_{r_j})$~\cite{p2,p6}. A step has a severity, and an option is
\emph{admissible} for the step when its severity-weighted worst-case expected loss, over the models
compatible with the evidence, fits a tolerated loss. Admissibility has a robust reading, for a
deployed build, and an ex-ante reading, for a build-and-evidence experiment with a statistical
certificate~\cite{p6}. On a directed acyclic dependency graph, the \emph{structural severity} of an
escape at node $v$ is $S(v)=\sum_w\rho(v\to w)\,c(w)\,\iota(w)$, where $\rho(v\to w)$ is the
probability that the fault reaches $w$, $c(w)$ the cost of reaching $w$, and $\iota(w)$ the
probability that this damage is not undone~\cite{p6}.

\paragraph{The allocation model.}
\begin{definition}[Allocation model]\label{def:model}
Unit classes $k$ have demand $\lambda_k\ge0$ per period. $\mathcal A_k$ is the finite set of options
admissible for $k$ under an explicitly chosen reading of admissibility~\cite{p6}. Option $a$ for class
$k$ consumes $D_{r,ak}\ge0$ of resource $r$ per unit and costs $c_{ak}$. Resource $r$ has capacity
$\kappa_r$ per period. The \emph{fluid model} is
\begin{equation}\label{eq:fluid}\tag{F}
x_{ak}\ge0,\qquad \sum_{a\in\mathcal A_k}x_{ak}=\lambda_k,\qquad \sum_{k,a}D_{r,ak}x_{ak}\le\kappa_r .
\end{equation}
Four qualifications belong to the definition.
\begin{enumerate}[label=(\arabic*),itemsep=2pt]
\item Admissibility must hold on the routed tasks and histories, including the chosen concurrency and
review arrangement. Benchmark averages do not establish it.
\item An ex-ante statistical certificate keeps that label; it does not become robust admissibility.
\item Attention time and CPU time enter as workload rates. Memory, simultaneous seats, precedence, and
deadlines need scheduling constraints beyond~\eqref{eq:fluid} (Proposition~\ref{prop:scheduling}).
\item Executor error laws may stay unknown within the bounds of~\cite{p6}. Arrival and service laws
need separate assumptions wherever a queueing result uses them.
\end{enumerate}
\end{definition}
We write $D_{ak}$ for the vector $(D_{r,ak})_r$ and $\kappa$ for $(\kappa_r)_r$.

\section{The allocation backbone}\label{sec:backbone}
\begin{theorem}[Admissible allocation is invariant to the executor's error law]\label{thm:invariance}
Suppose that bounds of the kind in~\cite{p6} decide every option's admissibility. Suppose also that
the coefficients $D$, $c$, $\lambda$, $\kappa$ do not depend on the error law, or are replaced by
bounds that hold for every law in the class. Then the feasible set and the optimal allocations
of~\eqref{eq:fluid} are the same for every executor error law, dependence structure, and routing
mixture that satisfies those bounds.
\end{theorem}
\begin{proof}
The program's data are $\lambda$, $\kappa$, $D$, $c$, and the sets $\mathcal A_k$. The bounds fix
each $\mathcal A_k$, by the admissibility theorem and the mixture lemma of~\cite{p6}. Equal data give
equal feasible and optimal sets.
\end{proof}
\begin{remark}
Retries consume attempts, so $D$ can depend on the error law. Then the coefficients need interval
data and the worst case. Admissibility must hold under the chosen concurrency
(Definition~\ref{def:model}(1)). Replacing law-dependent coefficients by bounds that hold for every
law is the worst-case reformulation of inexact linear programs~\cite{soyster73}; robust linear
programming generalizes it~\cite{bental99}.
\end{remark}

\begin{proposition}[Bottleneck bound; operational laws~\cite{buzen76,denning}]\label{prop:bottleneck}
If completed class-$k$ units require long-run average demand $D_{r,k}>0$ of resource $r$, then the
class-$k$ throughput satisfies $X_k\le\min_{r:D_{r,k}>0}\kappa_r/D_{r,k}$. Other classes' use of $r$
tightens the bound. If a fraction $f>0$ of units needs principal review with mean time $s$ given
review, then $Xfs\le\kappa_{\pr}$.
\end{proposition}
\begin{proof}
Over a horizon $T$, required work cannot exceed capacity, apart from boundary work of order $o(T)$.
Divide by $T$.
\end{proof}
\begin{remark}\label{rem:bottleneck}
With $\kappa_{\pr}=120$ principal minutes per period, $f=0.1$, and $s=5$ minutes, the ceiling is
$X\le240$ units per period; halving $f$ at the same $s$ doubles it to $480$. Admissible mechanization
does not imply that $f$ falls. It can keep every principal approval, add validation, or raise
exceptions. A lower $f$ with a longer $s$ can raise $fs$.
\end{remark}

\begin{theorem}[Feasibility of admissible routing for single-resource pools; Hall~\cite{hall},
Gale~\cite{gale57}, Ford and Fulkerson~\cite{fordfulkerson}]\label{thm:hall}
Suppose each option is a pool with capacity $b_a$, each unit uses one unit of its chosen pool, and no
other shared constraint exists. Write $\mathcal A(T)=\bigcup_{k\in T}\mathcal A_k$. A feasible routing
exists if and only if
\[
\sum_{k\in T}\lambda_k\le\sum_{a\in\mathcal A(T)}b_a\qquad\text{for every set }T\text{ of classes}.
\]
Integral data admit an integral routing. A minimum cut in the source--class--pool--sink network names
a violating set.
\end{theorem}
\begin{proof}
Use max-flow min-cut. A finite cut is fixed by the set $T$ of classes on the source side, and its
capacity is $\sum_{k\notin T}\lambda_k+\sum_{a\in\mathcal A(T)}b_a$. The routing exists exactly when
the maximum flow equals $\sum_k\lambda_k$. Integrality of the maximum flow gives the integral routing.
\end{proof}
The condition is the bipartite case of Gale's supply--demand feasibility theorem~\cite{gale57}, which
extends Hall's theorem.
\begin{example}[Shared gate]\label{ex:sharedgate}
The theorem is false when options share a resource. Take one class with demand $2$ and two options,
each a pool of capacity $1$, that both consume one principal gate of capacity $1$. The pool condition
holds ($2\le1+1$), but the gate can carry only one unit. Theorem~\ref{thm:bundle} covers this case.
\end{example}

\begin{theorem}[Exact feasibility certificate for resource bundles; Farkas~\cite{farkas}, LP
duality~\cite{bertsimas}]\label{thm:bundle}
Let every $\mathcal A_k$ be finite and nonempty. Then~\eqref{eq:fluid} is feasible if and only if, for
every price vector $p\ge0$,
\[
\sum_k\lambda_k\min_{a\in\mathcal A_k}p^\top D_{ak}\ \le\ p^\top\kappa .
\]
Infeasibility therefore has a certificate of nonnegative resource weights, not necessarily a subset
cut.
\end{theorem}
\begin{proof}
Necessity: weight the capacity rows by $p$ and use $\sum_ax_{ak}=\lambda_k$. Sufficiency: the dual of
the feasibility LP is
\[
\max\big\{y^\top\lambda-p^\top\kappa:\ y_k\le p^\top D_{ak}\ \forall a\in\mathcal A_k,\ p\ge0\big\}.
\]
The primal is infeasible if and only if the dual is unbounded, if and only if some $(y,p)$ has positive
value, if and only if the displayed inequality fails for some $p\ge0$.
\end{proof}
\begin{example}[Coordinatewise tests are not enough]\label{ex:coordinatewise}
One unit has options with demand vectors $(2,0)$ and $(0,2)$, and the capacities are $(0.6,0.6)$.
For each resource separately, some option uses none of it, so both coordinatewise tests pass. The price
$p=(1,1)$ gives $\min_a p^\top D_a=2>1.2=p^\top\kappa$ and certifies infeasibility. A certified
additive damage allowance can enter the model as one more resource row.
\end{example}
Suppose the resource rows are the edges of a capacitated network, each demand vector $D_{ak}$ is the
edge-incidence vector of a path, and each class's options are all paths between its source and its
sink. Then $\min_a p^\top D_{ak}$ is a shortest-path length under the edge lengths $p$, and
Theorem~\ref{thm:bundle} becomes the multicommodity feasibility theorem of Onaga and Kakusho and of
Iri~\cite{onaga71,iri71}. With only some paths admissible, the minimum is over the admissible paths.

\begin{theorem}[Shadow prices of scarce resources; LP duality and
sensitivity~\cite{boyd}]\label{thm:prices}
Let $V(\lambda,\kappa)=\min\{c^\top x:\ \eqref{eq:fluid}\}$, with dual
\[
\max\big\{y^\top\lambda-p^\top\kappa:\ y_k\le c_{ak}+p^\top D_{ak},\ p\ge0\big\}.
\]
\begin{enumerate}[label=(\roman*),itemsep=2pt]
\item At an optimal dual $(y,p)$, $V(\lambda,\kappa+\Delta)\ge V(\lambda,\kappa)-p^\top\Delta$. Where
$V$ is differentiable, $\partial V/\partial\kappa_r=-p_r$. At a kink, keep the set of optimal duals.
\item If the old problem is feasible, adding options (columns) strictly lowers the optimum if and only
if \emph{no} optimal dual of the old problem satisfies all the new columns' dual constraints. Without
feasibility the condition can hold vacuously. With a unique optimal dual, a negative reduced cost
$c+p^\top D-y_k<0$ is sufficient.
\item Changes to demand, coefficients, admissible sets, or available options are discrete and in
general need reoptimization. The prices give only local or bounding information about them.
\end{enumerate}
\end{theorem}
\begin{proof}
Part (i) is weak and strong duality. Part (ii): the new optimum equals the old one if and only if some
old optimal dual remains feasible for the enlarged dual. Part (iii) states the scope of (i) and (ii):
both are local or bounding statements.
\end{proof}
\begin{example}[The unrestricted reduced-cost test fails]\label{ex:degenerate}
Take the old problem $x=1$, $x\le1$, with cost $0$. Every $y=p\ge0$ is dual-optimal. A new option with
zero cost and zero resource use has reduced cost $0+0-1=-1$ at $y=p=1$, yet the optimum stays $0$.
\end{example}
The old problem is primal degenerate, so its optimal dual is not unique. That the prices of a single
dual solution can then mislead is known~\cite{aucamp82,akgul84,gal86,jansen97}, and directional
marginal values are extrema over the set of optimal duals~\cite{mills56,williams63}. Part~(ii) states
this fact for added columns.
The price $p_{\pr}$ values a principal minute at the current load and admissibility. A candidate
mechanization can then be compared across all scarce resources in one currency, with the limits of
parts (ii) and (iii).

\begin{proposition}[Few coupling constraints imply few split classes; basic feasible
solutions~\cite{bertsimas}, Lenstra, Shmoys, and Tardos~\cite{lst90}]\label{prop:split}
A feasible bounded instance of~\eqref{eq:fluid} with $K$ classes of positive demand and $m$ resource
rows has an optimal solution with at most $K+m$ positive coordinates. Hence at most $m$ classes are
split across options. At a vertex with $b$ binding independent resource rows, at most $b$ classes are
split.
\end{proposition}
\begin{proof}
A linear objective attains its optimum at a vertex. A vertex has at most as many positive coordinates
as independent active equality rows. Each class with positive demand needs at least one positive
coordinate, so at most $m$ (or $b$) classes have two or more.
\end{proof}
A fluid optimum therefore assigns most classes to one option. Randomized routing is needed for at
most as many classes as there are binding scarce resources. The same vertex count bounds the
fractional jobs in the linear relaxation of scheduling on unrelated parallel machines~\cite{lst90}, and
it gives at most three split categories in a recent model of AI review before human
approval~\cite{linzhang26}.

\section{The principal's attention and review sets}\label{sec:principal}
\begin{proposition}[Principal-attention allocation is a knapsack; Dantzig~\cite{dantzig},
Kellerer, Pferschy, and Pisinger~\cite{kellerer}]\label{prop:knapsack}
Suppose optional reviews have fixed avoided losses $v_i\ge0$ and times $t_i>0$, effects are additive
across units, and there are no setup or interaction effects. Then choosing reviews within principal
time $H$ is a knapsack problem.
\begin{enumerate}[label=(\alph*),itemsep=2pt]
\item The ratio order $v_i/t_i$ solves the fractional problem.
\item For the integer problem, first discard reviews with $t_i>H$. Then the better of the ratio prefix
of the remaining reviews and the best remaining single review is a $1/2$-approximation (the standard
modified greedy).
\item Without the discard step the guarantee fails. With $H=3$, one review $(t,v)=(4,100)$, and three
reviews $(1,1)$, an undiscarded ratio prefix is empty, and the best feasible single review gives $1$
against an optimum of $3$.
\item With complete blind sets, constant severity, and review that only suppresses faults,
$v_i=S_ip_i\mu_i(\Sigma_i\setminus\Sigma_{\pr})$. Here $p_i\mu_i$ is the unit's escape law through its
existing stack, $\Sigma_i$ is the blind set of that stack, $\Sigma_{\pr}$ is the principal's blind set,
and $S_i$ is the unit's severity.
\end{enumerate}
\end{proposition}
\begin{proof}
Parts (a) and (b) are the classical results~\cite{dantzig,kellerer}. Part (c): the review $(4,100)$
has the largest ratio and does not fit, so the prefix stops before it. Part~(c) shows why the
discard step, often stated as the assumption that every item fits, is needed. Part (d): the principal's review
removes the escaped faults outside $\Sigma_{\pr}$, each at damage $S_i$.
\end{proof}
\begin{remark}[Conditions]
Reviews required for admissibility are not optional and are scheduled first. Fault-dependent damage
needs integration, not a common severity. Lower bounds $\underline v_i$ give a conservative objective,
but a $1/2$-approximation to that objective is not a $1/2$-approximation to the actual benefit.
Overlapping coverage breaks additivity (Theorem~\ref{thm:submodular}).
\end{remark}

\begin{theorem}[Fixed-law review value is submodular; robust review value need not
be; Nemhauser, Wolsey, and Fisher~\cite{nemhauser}, Krause et al.~\cite{krause08}]\label{thm:submodular}
Fix a fault measure $\nu$, a damage $h\ge0$, and reviewer rejection sets $C_r$, where review only
suppresses faults. Then
\[
F_\nu(A)=\int h\,\ind\Big\{z\in\bigcup_{r\in A}C_r\Big\}\,d\nu(z)
\]
is monotone submodular, so greedy selection of $k$ reviewers of equal cost achieves a fraction $1-1/e$
of the optimum~\cite{nemhauser}. The robust value $F_{\mathrm{rob}}(A)=\inf_{\nu\in\mathcal U}F_\nu(A)$,
over a set $\mathcal U$ of fault measures, need not be submodular.
\end{theorem}
\begin{proof}
Fixed law: the marginal value of $r$ given a coalition is the damage-weighted mass that $r$ catches and
the coalition does not. This mass shrinks as the coalition grows. The $1-1/e$ guarantee is the theorem
of Nemhauser, Wolsey, and Fisher~\cite{nemhauser}. Robust counterexample: take faults $x,y$, unit
damage, $C_1=\{x\}$, $C_2=\{y\}$, and $\mathcal U=\{\delta_x,\delta_y\}$. Then
$F_{\mathrm{rob}}(\varnothing)=F_{\mathrm{rob}}(\{1\})=F_{\mathrm{rob}}(\{2\})=0$ and
$F_{\mathrm{rob}}(\{1,2\})=1$, so
$F_{\mathrm{rob}}(\{1\})+F_{\mathrm{rob}}(\{2\})<F_{\mathrm{rob}}(\{1,2\})+F_{\mathrm{rob}}(\varnothing)$.
\end{proof}
The fixed-law value is a weighted coverage function, and the $1-1/e$ guarantee is tight for
coverage~\cite{feige98}. The robust half is known: a pointwise minimum of monotone submodular
functions is in general not submodular, and robust selection is hard to approximate~\cite{krause08};
see also~\cite{orlin18}. Our counterexample has the same two-point form as theirs.
A reviewer can thus have zero guaranteed value alone and positive guaranteed value in a coalition.
Certified reviewer \emph{sets} should be allocated; reviewers should not be ranked by singleton lower
bounds.

\begin{corollary}[The principal as producer inherits common blind sets]\label{cor:blind}
Suppose the principal produces bad outputs with probability $p_{\pr}$ and conditional fault law
$\mu_{\pr}$, and every reviewer in an allowed coalition accepts a common set $B$. Then the
first-attempt escape floor of the principal's work is at least $p_{\pr}\mu_{\pr}(B)$. Self-review by
the principal, or more members that keep $B$, cannot lower it, provided the principal's own review
also accepts $B$. Production by the principal alone does not establish that premise.
\end{corollary}
\begin{proof}
Every fault in $B$ passes every member. Apply the conjunctive-review result of~\cite{p6} with the
principal as producer.
\end{proof}

\begin{proposition}[Implementation verification does not identify intent; Le~Cam's two-point
bound~\cite{lecam73}]\label{prop:intent}
Suppose two intended specifications induce the same available transcript but require opposite binary
actions, with loss $S$ for the wrong action. Then every transcript-based rule has worst-case expected
loss at least $S/2$. This holds even if the rule proves that its implementation satisfies the recorded
specification exactly.
\end{proposition}
\begin{proof}
A rule sees the same transcript under both specifications, so it chooses action~1 with the same
probability $q$. Its losses are $Sq$ and $S(1-q)$, and $\max\{Sq,S(1-q)\}\ge S/2$.
\end{proof}
This is the exact-twin case of the twin regret of~\cite{p6}, with equal gaps, and so an instance of
Le~Cam's two-point bound~\cite{lecam73}. The two specifications are observationally equivalent in
the sense of Koopmans~\cite{koopmans49}.
The split between building the right product and building the product right is Boehm's distinction
between validation and verification~\cite{boehm84}.
Establishing the operative specification (authority) and showing that an action meets it
(verification) are different objects. Before any theorem about the principal's errors, one must
define whether the principal's declaration \emph{is} the specification or reports a latent intent
imperfectly.

\section{Containment and organization}\label{sec:containment}
In this section, condition on an escape at node $v$. Let $S(v)$ be its expected damage, and
$S_{\mathrm{in}}(v)$ the expected damage inside $v$'s module.

\begin{theorem}[Modularity bounds severity]\label{thm:modularity}
\begin{enumerate}[label=(\alph*),itemsep=2pt]
\item \emph{General form.} Let $C$ be a set of interfaces that separates $v$'s module from the outside,
and let $E_e$ be the event that the fault's cascade ever crosses interface $e$. Suppose
$\Prb(E_e\mid\text{escape at }v,\text{routed history},\text{fault type})\le\varepsilon_e$ uniformly over
compatible models. Suppose also that outside damage is $0$ without a crossing and at most
$C_{\mathrm{out}}$ pathwise otherwise. Then
\[
S(v)\le S_{\mathrm{in}}(v)+C_{\mathrm{out}}\min\Big(1,\sum_{e\in C}\varepsilon_e\Big).
\]
No independence is needed.
\item \emph{Model form.} In the independent-edge model of~\cite{p6} on a directed acyclic graph, let $M$
be $v$'s module. Suppose every path from $v$ to a node outside $M$ uses an edge of $E_{\mathrm{cut}}$,
and each such edge transmits with probability at most $\varepsilon_e$. Write
$S_M(v)=\sum_{w\in M}\rho(v\to w)c(w)\iota(w)$ and $C_{\mathrm{out}}=\sum_{w\notin M}c(w)\iota(w)$. Then
\[
S(v)\le S_M(v)+C_{\mathrm{out}}\sum_{e\in E_{\mathrm{cut}}}\varepsilon_e\,\rho(v\to\mathrm{tail}(e)).
\]
\end{enumerate}
\end{theorem}
\begin{proof}
(a) Outside damage is at most $C_{\mathrm{out}}\ind_{\bigcup_eE_e}$. Take expectations and use the
union bound and $\Prb\le1$. (b) Reaching $w\notin M$ requires reaching $\mathrm{tail}(e)$ for some
$e\in E_{\mathrm{cut}}$ and $e$ transmitting. In a directed acyclic graph with independent edges these
two events are independent, because reaching $\mathrm{tail}(e)$ does not use $e$. Apply the union
bound, multiply by $c(w)\iota(w)$, and sum over $w\notin M$.
\end{proof}
\begin{remark}[Conditions]
A pass rate on a reference distribution is not a conditional crossing bound for a fault routed to the
interface. A per-attempt bound is not an ``ever crosses'' bound under retries. Contracts must
intercept the relevant semantic faults, be enforced, and have no bypass. Smaller severity weakly
enlarges the budget of~\cite{p6}; it does not guarantee a cheaper admissible option, or that the
interface overhead pays.
\end{remark}
Part~(a) is the union-bound form of the cut bounds of reliability theory~\cite{esary63}, without
their independence assumption. The ceilings $\varepsilon_e$ play the role of the error permeabilities
of software modules~\cite{hiller01,hiller04}, and the module bounded by $C$ is a fault-containment
region~\cite{lala94}.

\begin{theorem}[Serial containment multiplies conditional leakage ceilings]\label{thm:serial}
Suppose every route from an escaped fault to outside damage must survive ordered stages $1,\dots,L$.
Let $B_0$ be the escape and $B_j\subseteq B_{j-1}$ the event of surviving through stage $j$, and
suppose $\Prb(B_j\mid B_{j-1},\text{history entering }j)\le\varepsilon_j$ uniformly. Then
$\Prb(B_L)\le\prod_j\varepsilon_j$, and, with $C_{\mathrm{out}}$ as in
Theorem~\ref{thm:modularity}(a),
\[
S(v)\le S_{\mathrm{in}}(v)+C_{\mathrm{out}}\prod_{j=1}^L\varepsilon_j .
\]
\end{theorem}
\begin{proof}
$\Prb(B_j)\le\varepsilon_j\Prb(B_{j-1})$; iterate. Outside damage is at most
$C_{\mathrm{out}}\ind_{B_L}$; take expectations.
\end{proof}
Parallel escape routes add (Theorem~\ref{thm:modularity}, union bound). Serial stages multiply
(Theorem~\ref{thm:serial}, conditional bounds). Marginal ceilings do not multiply: two stages with
marginal failure probability $\varepsilon$ but identical failure events give joint failure
$\varepsilon$, not $\varepsilon^2$. Renaming reviewers as stages buys no conditional guarantee.
For independent layers, the product is the layer-multiplication rule of layer-of-protection
analysis~\cite{lopa}; Theorem~\ref{thm:serial} replaces independence by uniform conditional ceilings.
The failure of marginal multiplication is the coincident-failure effect of diverse
software~\cite{eckhardt85,littlewood89}.

\begin{proposition}[Parallelism on coupled work; universal scalability law~\cite{gunther}]\label{prop:usl}
\begin{enumerate}[label=(\alph*),itemsep=2pt]
\item If participating workers are optional, so that extra workers may stay idle, the best achievable
throughput is nondecreasing in the number of available workers.
\item Suppose throughput follows $X(n)=X(1)\,n/[1+\alpha(n-1)+\beta n(n-1)]$ with $0\le\alpha<1$ and
$\beta>0$, and $X(1)>0$. Its continuous maximum is at $n^\ast=\max\big(1,\sqrt{(1-\alpha)/\beta}\big)$. The law
is not globally concave; its tail is convex.
\item With $\beta=0$ and $X(1)>0$ there is no finite maximizer. For $0<\alpha<1$, throughput increases to the limit
$X(1)/\alpha$; for $\alpha=0$, it grows linearly without bound.
\end{enumerate}
\end{proposition}
\begin{proof}
(a) Idling the extra workers reproduces any smaller configuration. (b) The sign of the derivative is
the sign of $1-\alpha-\beta n^2$. $X$ is rational, so its Laurent expansion at infinity converges and may
be differentiated term by term: $X(n)=X(1)/(\beta n)-X(1)(\alpha-\beta)/(\beta^2n^2)+O(n^{-3})$, hence
$X''(n)=2X(1)/(\beta n^3)+O(n^{-4})>0$ for all sufficiently large $n$. (c)
With $\beta=0$, $X(n)=X(1)n/[1-\alpha+\alpha n]$, which increases in $n$ and tends to $X(1)/\alpha$ for
$\alpha>0$ and equals $X(1)n$ for $\alpha=0$.
\end{proof}
Part (a) is proven here. The law in (b) and (c) is Gunther's model~\cite{gunther}, and Gunther also
gives its maximizer and its reduction to Amdahl's law at $\beta=0$~\cite{gunther08}; Amdahl stated
that ceiling in words~\cite{amdahl67}. For example, $\alpha=0.1$ and $\beta=0.01$ give $n^\ast=\sqrt{90}\approx9.49$, and the
integer worker counts $9$ and $10$ tie, both giving $X(n)/X(1)=25/7$. A falling throughput curve describes a policy that makes
extra workers participate. It is not a limit on what is achievable. Pairwise coordination cost does
not by itself fix the law.

\begin{corollary}[Review lineages as capacity]\label{cor:lineages}
Treat each certified review bundle as an option with a resource-demand vector. Its throughput obeys
Proposition~\ref{prop:bottleneck} and Theorem~\ref{thm:bundle}, and Theorem~\ref{thm:hall} when the
pools are interchangeable single resources. Distinct lineages do not by themselves give independence or
complementary blind sets. What matters is the certified acceptance intersection for the producer and
the routed workload~\cite{p6}.
\end{corollary}
\begin{proof}
A bundle with a demand vector is an option of Definition~\ref{def:model}; apply the cited results. The
escape of a conjunctive stack depends on the intersection of its members' acceptance sets, not on
their lineage labels~\cite{p6}.
\end{proof}

\section{Substitution, scheduling, and stability}\label{sec:scheduling}
\begin{proposition}[Lifecycle substitution needs vector, time-indexed feasibility]\label{prop:lifecycle}
Over a horizon, let $s_r$ be the current slack of resource $r$, with the rest of the allocation fixed.
Replace $n$ units by a mechanization with fixed demand $F_r$ (build, validation) and recurring demand
$d^{\mathrm{cl}}_r$ in place of $d^{\mathrm{op}}_r$. The replacement is feasible if and only if
\[
F_r+n\big(d^{\mathrm{cl}}_r-d^{\mathrm{op}}_r\big)\le s_r\qquad\text{for every }r,
\]
and in every period when demands are time-indexed.
\end{proposition}
\begin{proof}
Substitute the changed demand into each capacity row.
\end{proof}
Multi-period capacity planning keeps a capacity row for each resource and each
period~\cite{luss82}.
\begin{example}[Lifecycle totals mislead]\label{ex:lifecycle}
Let the principal's capacity be $10$ in each of two periods, the old demand $(10,10)$, and the new
demand $(15,0)$. Total demand falls from $20$ to $15$, but period one is infeasible. A slack resource
can also become the new bottleneck. Reallocation elsewhere can compensate. Fixed costs and integer
commitments lie outside the marginal test of Theorem~\ref{thm:prices}.
\end{example}

\begin{proposition}[Queueing at a scarce gate; Harchol-Balter~\cite{harchol}, Kingman~\cite{kingman61,kingman62},
Cox and Smith~\cite{coxsmith}]\label{prop:queue}
Each item holds under its stated model.
\begin{enumerate}[label=(\arabic*),itemsep=2pt]
\item \emph{M/M/1, first come first served.} With Poisson arrivals at rate $\lambda$, exponential
service at rate $\mu>\lambda$, and one server, $\E W=1/(\mu-\lambda)$ and $\E W_q=\rho/(\mu-\lambda)$,
with $\rho=\lambda/\mu$.
\item \emph{Kingman (GI/GI/1).} $\E W_q\approx\frac{\rho}{1-\rho}\cdot\frac{c_a^2+c_s^2}{2}\,\E S$, where
$c_a$ and $c_s$ are the coefficients of variation of interarrival and service times. This is a
heavy-traffic approximation, not a universal bound.
\item \emph{The $c\mu$ rule.} With Poisson class arrivals, exponential class service, one
preemptive-resume server, no switching cost, and constant linear holding costs $c_k$, serving the
largest $c_k\mu_k$ first is optimal.
\item \emph{Batching.} With setup time $h$, batch size $b$, and per-unit service time $s$, saturated
capacity is $b/(h+bs)$. Under Poisson arrivals, forming a batch adds mean per-unit delay
$(b-1)/(2\lambda)$.
\end{enumerate}
\end{proposition}
\begin{proof}
Items (1)--(3) are classical~\cite{harchol,kingman61,kingman62,coxsmith,buyukkoc85}. Item (4): a saturated cycle
processes $b$ units in time $h+bs$. In a batch, the $i$-th arrival waits for $b-i$ further arrivals,
each after a mean interarrival time $1/\lambda$; the average of $(b-i)/\lambda$ over $i=1,\dots,b$ is
$(b-1)/(2\lambda)$. Batch-service queues go back to Bailey and Neuts~\cite{bailey54,neuts67}; item~(4)
gives only the elementary capacity and formation-delay terms.
\end{proof}
\begin{remark}[Caveats]
Severity is not automatically a holding-cost rate. Work awaiting safe release incurs delay cost, while
work already exposed accumulates damage. The invariance of~\cite{p6} to error laws gives no invariance
to arrival or service laws.
\end{remark}

\begin{proposition}[Verification compute is resource-constrained scheduling; Pritsker, Watters, and
Wolfe~\cite{pritsker69}]\label{prop:scheduling}
Suppose durations, machine eligibility, memory needs, releases, and deadlines are known. Then
allocating proof jobs is a cumulative-resource scheduling problem. Let $z_{jmt}\in\{0,1\}$ indicate that
job $j$ starts on machine $m$ at slot $t$, let $p_{jm}$ be its duration on $m$, $m_j$ its memory
need, and $M_m$ the memory of machine $m$. Memory requires
\[
\sum_{j,t:\ t\le\tau<t+p_{jm}}m_jz_{jmt}\le M_m\qquad\text{for every }m\text{ and }\tau,
\]
together with assignment, precedence, and deadline constraints. A cloud option adds placements and
costs. The LP relaxation gives lower bounds and prices in the sense of Theorem~\ref{thm:prices}.
\end{proposition}
\begin{proof}
A job started at slot $t$ on machine $m$ occupies $m_j$ units of memory in the slots
$t,\dots,t+p_{jm}-1$, so the displayed sum is the memory in use at slot $\tau$. Relaxing integrality
enlarges the feasible set, so the relaxation's optimal cost is a lower bound, and its dual solutions are
prices.
\end{proof}
\begin{example}[Memory-time totals are only a relaxation]\label{ex:memory}
Three one-slot jobs each need $6$ units of memory on a $10$-unit machine. No two of them can overlap,
since $2\cdot6=12>10$. Yet a rate of $1/2$ for each job passes the aggregate test
$3\cdot\tfrac12\cdot6=9\le10$, although a total rate of $3/2>1$ requires overlap in some slot.
\end{example}
The time-indexed formulation is that of Pritsker, Watters, and Wolfe~\cite{pritsker69}. Memory is a
cumulative resource~\cite{aggoun93}, and the aggregate test is an energy-type relaxation of the kind
studied in cumulative scheduling~\cite{baptiste01}.
\begin{remark}[Caveats]
Covered severity is not additive when proof scopes overlap. A solver run is an admissible verification
option only under its soundness, artifact-binding, and enforcement premises. A rule that all solver runs
go to the cloud is a policy constraint, not an optimality result.
\end{remark}

\begin{theorem}[MaxWeight scheduling of feasible configurations is throughput-optimal; Tassiulas and
Ephremides~\cite{tassiulas}]\label{thm:maxweight}
Consider slotted queues $Q_k(t+1)=[Q_k(t)-s_k(t)]_++A_k(t)$. Let $\mathcal M$ be a finite set of service
vectors that are simultaneously feasible, resource compatibility included. Let the arrivals be
i.i.d.\ with mean $\lambda$ and finite second moments, independent of a finite initial backlog with
$\E\|Q(0)\|_2^2<\infty$. If some $\bar s\in\conv\mathcal M$ has $\bar s\ge\lambda+\epsilon\ind$ with
$\epsilon>0$, then MaxWeight, $s(t)\in\arg\max_{s\in\mathcal M}Q(t)^\top s$, is strongly stable, with
\[
\limsup_{T\to\infty}\frac1T\sum_{t<T}\E\|Q(t)\|_1\le\frac{B}{2\epsilon},\qquad
B=\E\|A\|_2^2+\max_{s\in\mathcal M}\|s\|_2^2 .
\]
\end{theorem}
\begin{proof}[Proof sketch]
Use quadratic Lyapunov drift. One slot gives
$\E[\|Q(t+1)\|_2^2-\|Q(t)\|_2^2\mid Q(t)]\le B+2Q(t)^\top(\lambda-s(t))$. MaxWeight maximizes
$Q^\top s$ over $\mathcal M$, hence over $\conv\mathcal M$, so $Q^\top s(t)\ge Q^\top\bar s\ge
Q^\top\lambda+\epsilon\|Q\|_1$. The drift is therefore at most $B-2\epsilon\|Q(t)\|_1$. Take
expectations, telescope over $t<T$, use $\E\|Q(0)\|_2^2<\infty$, and divide by $2\epsilon T$.
The explicit bound is the standard Lyapunov-drift form~\cite{georgiadis06,neely10}.
\end{proof}
Build $\mathcal M$ from options whose safety bounds hold under each mode's concurrency. Rework that
depends on errors belongs in the arrival and service model. Stability does not enforce the lifetime
consequence budget of~\cite{p6}. Throughput sustained indefinitely at positive risk and a finite
lifetime budget can be incompatible.

\section{Consequences for scheduler design}\label{sec:design}
The results suggest the following guidance. It is design guidance, not law. Each item is only as strong
as the result it cites, and holds only under that result's conditions.
\begin{enumerate}[leftmargin=*,itemsep=3pt]
\item \emph{Filter first, then allocate.} Admit an option to $\mathcal A_k$ only if its admissibility
holds under the allocation that will run it: the routed tasks, the concurrency, and the review
arrangement (Definition~\ref{def:model}). Then the allocation does not need the executors' error laws
(Theorem~\ref{thm:invariance}). If retries make demand depend on the error law, plan against interval
data and the worst case.
\item \emph{Check feasibility jointly.} Test feasibility with a price certificate over all resources
(Theorem~\ref{thm:bundle}). Resource-by-resource tests can pass an infeasible plan
(Example~\ref{ex:coordinatewise}). Subset cuts are exact for single-resource pools but insufficient for general resource bundles
(Theorem~\ref{thm:hall}, Example~\ref{ex:sharedgate}).
\item \emph{Use prices locally; reoptimize for discrete changes.} Dual prices value scarce resources,
including principal minutes, in one currency. To decide whether a new option helps, check all old
optimal duals, or reoptimize (Theorem~\ref{thm:prices}, Example~\ref{ex:degenerate}).
\item \emph{Expect few split classes.} A fluid optimum splits at most as many classes as there are
binding resources (Proposition~\ref{prop:split}). Randomized routing can be confined to those classes.
\item \emph{Schedule required principal reviews first.} Then fill the remaining principal time with
optional reviews by the ratio rule, after discarding reviews that do not fit, and compare with the best
single review (Proposition~\ref{prop:knapsack}). A guarantee on lower-bound values is not a guarantee on
actual benefit.
\item \emph{Allocate certified reviewer sets.} Do not rank reviewers by singleton lower bounds; a
reviewer with no guaranteed value alone can complete a coalition (Theorem~\ref{thm:submodular}).
Lineage labels do not certify complementary blind sets (Corollary~\ref{cor:lineages}).
\item \emph{Keep authority and verification apart.} A proof that work meets the recorded
specification does not show that the specification is the intended one
(Proposition~\ref{prop:intent}). Work the principal produces keeps the common blind set of its
reviewers (Corollary~\ref{cor:blind}).
\item \emph{Design containment with conditional bounds.} Parallel interfaces add their crossing
ceilings; serial stages multiply only conditional ceilings (Theorems~\ref{thm:modularity}
and~\ref{thm:serial}). A benchmark pass rate is not a crossing bound for routed faults.
\item \emph{Let workers idle.} With optional participation, adding workers cannot lower the best
achievable throughput (Proposition~\ref{prop:usl}). A falling curve indicates a participation policy
to change.
\item \emph{Test substitutions per resource and per period.} A mechanization that lowers total demand
can still overload one period or one resource (Proposition~\ref{prop:lifecycle},
Example~\ref{ex:lifecycle}).
\item \emph{State queueing assumptions.} Queueing formulas need arrival and service assumptions that
the error-law invariance does not supply. Do not use severity as a holding-cost rate without a reason
(Proposition~\ref{prop:queue}).
\item \emph{Schedule memory, not memory-time.} Verification jobs with memory footprints need
cumulative scheduling constraints; aggregate memory-time is a relaxation
(Proposition~\ref{prop:scheduling}, Example~\ref{ex:memory}).
\item \emph{Stabilize over admissible configurations, and budget separately.} MaxWeight over
configurations built from admissible options keeps queues stable when demand lies strictly inside the
capacity region. It does not enforce a lifetime risk budget (Theorem~\ref{thm:maxweight}).
\end{enumerate}

\section{Related work}\label{sec:related}
Allocating inspection effort is an old problem. Lindsay and Bishop place screening inspections on a
serial production line by dynamic programming~\cite{lindsay64}, and surveys cover the
field~\cite{raz86,mandroli06}. Testing effort has been allocated across software modules under a
reliability-growth law~\cite{ohtera90}. These models assume a known fault law; our admissible sets use
only bounds. Chance-constrained programming requires a constraint to hold with a stated probability
under the decision taken~\cite{charnes59}. Our admissibility is a worst-case premise, and it must hold
under the chosen concurrency and review arrangement.

For AI systems, learning to defer routes cases from a model to a human~\cite{madras18,mozannar20}.
Some of this work adds an explicit limit on human workload~\cite{okati21,alves24} or congestion at the
human queue~\cite{lykouris24}. Lin and Zhang model agent-generated code changes that may pass an AI
reviewer before a human maintainer approves them, as a queueing network with three resource
pools~\cite{linzhang26}. They plan with a fluid linear program, use its dual prices, show that the
review policy changes at staffing breakpoints and cannot be produced by one fixed ranking, count at
most three split categories at an extreme point, as in Proposition~\ref{prop:split}, and prove a policy
asymptotically optimal. Their model is the closest to Sections~\ref{sec:backbone}
and~\ref{sec:scheduling}. It takes the false-acceptance and false-rejection rates as parameters. We first filter
options by admissibility from bounds; when the remaining program data are law-independent or replaced
by common uniform bounds, the allocation then does not depend on the error law
(Theorem~\ref{thm:invariance}). We also add the review-set and containment results. AI control spends a
limited human audit budget against a possibly adversarial model~\cite{greenblatt23,griffin24}, and
language-model routers trade cost against quality for each query~\cite{chen23,ong24}. Drift-plus-penalty
methods extend MaxWeight with time-average constraints~\cite{neely10}; they are a possible route to
the lifetime budget that Theorem~\ref{thm:maxweight} does not enforce.

Requiring feasibility under the decision actually chosen is itself standard: stochastic and robust
optimization with decision-dependent uncertainty evaluate constraints against the uncertainty that the
chosen decision induces~\cite{goel06,nohadani18}. What we add is narrow: the application of that
requirement to gated work, where the uncertain premise is admissibility under the chosen concurrency
and review arrangement; the conditions under which the classical tools cited above then apply; and
counterexamples where they fail.

\section{Limitations}\label{sec:limits}
The proofs are by hand; no proof assistant checked them. The finite checks of
Section~\ref{sec:repro} test instances and counterexamples; they do not replace the proofs. The fluid
model is a model: executable schedules need the constraints of Proposition~\ref{prop:scheduling}
and Theorem~\ref{thm:maxweight}. Real workloads are not shown to satisfy any queueing or arrival
assumption used here. Admissibility enters as an input per class under a stated reading; it must hold
under the chosen concurrency and routing, and this paper does not measure it. The results do not
estimate any capacity, price, coverage, or severity for a real system.

Two directions are not addressed here. The first is measurement design: how to allocate a limited audit budget
to estimate the bounds that the scheduler relies on. Information lower bounds for adaptive sampling are
a starting point~\cite{kaufmann}. The second is feedback: a policy decision changes constraints, and
the allocation changes demand in turn. Contraction arguments give a unique equilibrium and a Lipschitz
sensitivity when they apply~\cite{banach}, but LP-based allocations can jump at ties, so continuity
must be established, not assumed. The principal as an executor is treated here only through
Corollary~\ref{cor:blind} and Proposition~\ref{prop:intent}; a model of intent elicitation is needed
before more can be said. This is a preprint. It has not been peer reviewed. The prior-art search was targeted, not
systematic.

\section{Reproducibility}\label{sec:repro}
The script \texttt{simulate\_allocation.py} reproduces every numerical claim and counterexample in
this paper with exact rational arithmetic and the Python standard library. It prints PASS or FAIL for
each claim and exits with a nonzero status on any failure. It also checks random instances: Hall's
condition against maximum flow, the price certificate against direct feasibility, strong duality and
the reduced-cost test on a small LP, the fractional and integer knapsack guarantees, submodularity and
the greedy guarantee for fixed-law review value, the modularity bound by exact enumeration on a small
graph, and the MaxWeight drift inequality. Four negative controls assert false claims (the shared
gate, the integer ratio rule, the degenerate dual, and serial stages from marginals) and must be
rejected.

\begin{center}\small
\begin{tabular}{p{0.66\textwidth}r}
\toprule Quantity & Value\\\midrule
Remark~\ref{rem:bottleneck}: principal-review ceiling, $f=0.1$ / $f=0.05$ & 240 / 480\\
Example~\ref{ex:coordinatewise}: $\min_a p^\top D_a$ versus $p^\top\kappa$ at $p=(1,1)$ & 2 versus 1.2\\
Example~\ref{ex:degenerate}: reduced cost at $y=p=1$; old and new optimum & $-1$; 0 and 0\\
Proposition~\ref{prop:knapsack}(c): undiscarded greedy / optimum & 1 / 3\\
Theorem~\ref{thm:submodular}: robust values of $\varnothing,\{1\},\{2\},\{1,2\}$ & 0, 0, 0, 1\\
Proposition~\ref{prop:usl}: $n^\ast$ for $\alpha=0.1$, $\beta=0.01$ & $\sqrt{90}\approx9.49$\\
Example~\ref{ex:lifecycle}: total demand old / new; period-one demand & 20 / 15; 15 $>$ 10\\
Example~\ref{ex:memory}: aggregate memory-time / capacity & 9 / 10\\
\bottomrule
\end{tabular}
\end{center}

\section*{Code and data availability}
The standard-library Python script \texttt{simulate\_allocation.py} reproduces every numerical claim of
Section~\ref{sec:repro}. It is at \url{https://github.com/ivmat/gated-computation-sim}.
This paper is archived on Zenodo, DOI~\href{https://doi.org/10.5281/zenodo.23449199}{10.5281/zenodo.23449199}.

\begin{thebibliography}{99}
\bibitem{p1} I. Matija\v{s}evi\'c. A Fault-Tolerance Threshold for Gated Agentic Computation.
2026. \url{https://doi.org/10.5281/zenodo.20820968}.
\bibitem{p2} I. Matija\v{s}evi\'c. Generated Gates Inherit Their Generator's Blind Spots.
2026. \url{https://doi.org/10.5281/zenodo.20837102}.
\bibitem{p3} I. Matija\v{s}evi\'c. Escape, Cost, and Correlation at the Verification Floor of
Gated Agentic Computation. 2026. \url{https://doi.org/10.5281/zenodo.21456783}.
\bibitem{p4} I. Matija\v{s}evi\'c. Progressive Semantic Mechanization. 2026.
\url{https://doi.org/10.5281/zenodo.23119225}.
\bibitem{p5} I. Matija\v{s}evi\'c. Beyond Constant Fault Rates: Reliability Bounds for Gated Computation and Semantic Mechanization. 
Preprint, 2026. \url{https://doi.org/10.5281/zenodo.23293740}.
\bibitem{p6} I. Matija\v{s}evi\'c. Risk-Gated Mechanization: Severity, Verifiability, and Decisions Invariant to the Executor's Error Law. 
Preprint, 2026. \url{https://doi.org/10.5281/zenodo.23324667}.
\bibitem{aggoun93} A. Aggoun and N. Beldiceanu. Extending CHIP in order to solve complex scheduling
and placement problems. \emph{Mathematical and Computer Modelling} 17(7):57--73, 1993.
\bibitem{akgul84} M. Akg\"ul. A note on shadow prices in linear programming. \emph{Journal of the
Operational Research Society} 35(5):425--431, 1984.
\bibitem{alves24} J. V. Alves, D. Leit\~ao, S. Jesus, M. O. P. Sampaio, J. Li\'ebana, P. Saleiro,
M. A. T. Figueiredo, and P. Bizarro. Cost-sensitive learning to defer to multiple experts with workload
constraints. 2024. \url{https://arxiv.org/abs/2403.06906}.
\bibitem{amdahl67} G. M. Amdahl. Validity of the single processor approach to achieving large scale
computing capabilities. \emph{AFIPS Spring Joint Computer Conference}, pp.~483--485, 1967.
\bibitem{aucamp82} D. C. Aucamp and D. I. Steinberg. The computation of shadow prices in linear
programming. \emph{Journal of the Operational Research Society} 33(6):557--565, 1982.
\bibitem{bailey54} N. T. J. Bailey. On queueing processes with bulk service. \emph{Journal of the
Royal Statistical Society, Series B} 16(1):80--87, 1954.
\bibitem{banach} S. Banach. Sur les op\'erations dans les ensembles abstraits et leur application aux
\'equations int\'egrales. \emph{Fundamenta Mathematicae} 3:133--181, 1922.
\bibitem{baptiste01} P. Baptiste, C. Le~Pape, and W. Nuijten. \emph{Constraint-Based Scheduling}.
Kluwer, 2001.
\bibitem{bental99} A. Ben-Tal and A. Nemirovski. Robust solutions of uncertain linear programs.
\emph{Operations Research Letters} 25(1):1--13, 1999.
\bibitem{bertsimas} D. Bertsimas and J. N. Tsitsiklis. \emph{Introduction to Linear Optimization}.
Athena Scientific, 1997.
\bibitem{boehm84} B. W. Boehm. Verifying and validating software requirements and design
specifications. \emph{IEEE Software} 1(1):75--88, 1984.
\bibitem{boyd} S. Boyd and L. Vandenberghe. \emph{Convex Optimization}. Cambridge University Press,
2004.
\bibitem{buyukkoc85} C. Buyukkoc, P. Varaiya, and J. Walrand. The $c\mu$ rule revisited. \emph{Advances
in Applied Probability} 17(1):237--238, 1985.
\bibitem{buzen76} J. P. Buzen. Fundamental operational laws of computer system performance.
\emph{Acta Informatica} 7(2):167--182, 1976.
\bibitem{charnes59} A. Charnes and W. W. Cooper. Chance-constrained programming. \emph{Management
Science} 6(1):73--79, 1959.
\bibitem{chen23} L. Chen, M. Zaharia, and J. Zou. FrugalGPT: how to use large language models while
reducing cost and improving performance. 2023. \url{https://arxiv.org/abs/2305.05176}.
\bibitem{coxsmith} D. R. Cox and W. L. Smith. \emph{Queues}. Methuen, 1961.
\bibitem{dantzig} G. B. Dantzig. Discrete-variable extremum problems. \emph{Operations Research}
5(2):266--288, 1957.
\bibitem{denning} P. J. Denning and J. P. Buzen. The operational analysis of queueing network models.
\emph{ACM Computing Surveys} 10(3):225--261, 1978.
\bibitem{eckhardt85} D. E. Eckhardt and L. D. Lee. A theoretical basis for the analysis of multiversion
software subject to coincident errors. \emph{IEEE Transactions on Software Engineering}
SE-11(12):1511--1517, 1985.
\bibitem{esary63} J. D. Esary and F. Proschan. Coherent structures of non-identical components.
\emph{Technometrics} 5(2):191--209, 1963.
\bibitem{farkas} J. Farkas. Theorie der einfachen Ungleichungen. \emph{Journal f\"ur die reine und
angewandte Mathematik} 124:1--27, 1902.
\bibitem{feige98} U. Feige. A threshold of $\ln n$ for approximating set cover. \emph{Journal of the
ACM} 45(4):634--652, 1998.
\bibitem{fordfulkerson} L. R. Ford and D. R. Fulkerson. Maximal flow through a network.
\emph{Canadian Journal of Mathematics} 8:399--404, 1956.
\bibitem{goel06} V. Goel and I. E. Grossmann. A class of stochastic programs with decision dependent
uncertainty. \emph{Mathematical Programming} 108:355--394, 2006.
\bibitem{gal86} T. Gal. Shadow prices and sensitivity analysis in linear programming under degeneracy.
\emph{OR Spektrum} 8(2):59--71, 1986.
\bibitem{gale57} D. Gale. A theorem on flows in networks. \emph{Pacific Journal of Mathematics}
7(2):1073--1082, 1957.
\bibitem{georgiadis06} L. Georgiadis, M. J. Neely, and L. Tassiulas. Resource allocation and
cross-layer control in wireless networks. \emph{Foundations and Trends in Networking} 1(1):1--144,
2006.
\bibitem{griffin24} C. Griffin, L. Thomson, B. Shlegeris, and A. Abate. Games for AI control: models of
safety evaluations of AI deployment protocols. 2024. \url{https://arxiv.org/abs/2409.07985}.
\bibitem{greenblatt23} R. Greenblatt, B. Shlegeris, K. Sachan, and F. Roger. AI control: improving
safety despite intentional subversion. 2023. \url{https://arxiv.org/abs/2312.06942}.
\bibitem{gunther} N. J. Gunther. \emph{Guerrilla Capacity Planning}. Springer, 2007.
\bibitem{gunther08} N. J. Gunther. A general theory of computational scalability based on rational
functions. 2008. \url{https://arxiv.org/abs/0808.1431}.
\bibitem{hall} P. Hall. On representatives of subsets. \emph{Journal of the London Mathematical
Society} 10:26--30, 1935.
\bibitem{harchol} M. Harchol-Balter. \emph{Performance Modeling and Design of Computer Systems:
Queueing Theory in Action}. Cambridge University Press, 2013.
\bibitem{hiller01} M. Hiller, A. Jhumka, and N. Suri. An approach for analysing the propagation of
data errors in software. \emph{Proceedings of the International Conference on Dependable Systems and
Networks}, pp.~161--170, 2001.
\bibitem{hiller04} M. Hiller, A. Jhumka, and N. Suri. EPIC: profiling the propagation and effect of
data errors in software. \emph{IEEE Transactions on Computers} 53(5):512--530, 2004.
\bibitem{iri71} M. Iri. On an extension of the maximum-flow minimum-cut theorem to multicommodity
flows. \emph{Journal of the Operations Research Society of Japan} 13(3):129--135, 1971.
\bibitem{jansen97} B. Jansen, J. J. de~Jong, C. Roos, and T. Terlaky. Sensitivity analysis in linear
programming: just be careful! \emph{European Journal of Operational Research} 101(1):15--28, 1997.
\bibitem{kaufmann} E. Kaufmann, O. Capp\'e, and A. Garivier. On the complexity of best-arm
identification in multi-armed bandit models. \emph{Journal of Machine Learning Research} 17:1--42,
2016.
\bibitem{kellerer} H. Kellerer, U. Pferschy, and D. Pisinger. \emph{Knapsack Problems}. Springer,
2004.
\bibitem{kingman61} J. F. C. Kingman. The single server queue in heavy traffic. \emph{Mathematical
Proceedings of the Cambridge Philosophical Society} 57(4):902--904, 1961.
\bibitem{kingman62} J. F. C. Kingman. On queues in heavy traffic. \emph{Journal of the Royal
Statistical Society, Series B} 24(2):383--392, 1962.
\bibitem{koopmans49} T. C. Koopmans. Identification problems in economic model construction.
\emph{Econometrica} 17(2):125--144, 1949.
\bibitem{krause08} A. Krause, H. B. McMahan, C. Guestrin, and A. Gupta. Robust submodular observation
selection. \emph{Journal of Machine Learning Research} 9:2761--2801, 2008.
\bibitem{lala94} J. H. Lala and R. E. Harper. Architectural principles for safety-critical real-time
applications. \emph{Proceedings of the IEEE} 82(1):25--40, 1994.
\bibitem{lecam73} L. Le~Cam. Convergence of estimates under dimensionality restrictions. \emph{Annals
of Statistics} 1(1):38--53, 1973.
\bibitem{lindsay64} G. F. Lindsay and A. B. Bishop. Allocation of screening inspection effort---a
dynamic-programming approach. \emph{Management Science} 10(2):342--352, 1964.
\bibitem{linzhang26} R. Lin and J. Zhang. When should AI review precede human approval of
agent-generated code? 2026, version~2; version~1 was titled ``When to screen, when to bypass:
LLM-judges in resource-scarce AI-human workflow''. \url{https://arxiv.org/abs/2603.13870}.
\bibitem{littlewood89} B. Littlewood and D. R. Miller. Conceptual modeling of coincident failures in
multiversion software. \emph{IEEE Transactions on Software Engineering} 15(12):1596--1614, 1989.
\bibitem{lopa} Center for Chemical Process Safety. \emph{Layer of Protection Analysis: Simplified
Process Risk Assessment}. American Institute of Chemical Engineers, 2001.
\bibitem{lst90} J. K. Lenstra, D. B. Shmoys, and \'E. Tardos. Approximation algorithms for scheduling
unrelated parallel machines. \emph{Mathematical Programming} 46:259--271, 1990.
\bibitem{luss82} H. Luss. Operations research and capacity expansion problems: a survey.
\emph{Operations Research} 30(5):907--947, 1982.
\bibitem{lykouris24} T. Lykouris and W. Weng. Learning to defer in congested systems: the AI-human
interplay. 2024. \url{https://arxiv.org/abs/2402.12237}.
\bibitem{madras18} D. Madras, T. Pitassi, and R. Zemel. Predict responsibly: improving fairness and
accuracy by learning to defer. 2017. \url{https://arxiv.org/abs/1711.06664}.
\bibitem{mandroli06} S. S. Mandroli, A. K. Shrivastava, and Y. Ding. A survey of inspection strategy
and sensor distribution studies in discrete-part manufacturing processes. \emph{IIE Transactions}
38(4):309--328, 2006.
\bibitem{mills56} H. D. Mills. Marginal values of matrix games and linear programs. In H. W. Kuhn and
A. W. Tucker (eds.), \emph{Linear Inequalities and Related Systems}, Annals of Mathematics Studies 38,
pp.~183--194. Princeton University Press, 1956.
\bibitem{mozannar20} H. Mozannar and D. Sontag. Consistent estimators for learning to defer to an
expert. 2020. \url{https://arxiv.org/abs/2006.01862}.
\bibitem{neely10} M. J. Neely. \emph{Stochastic Network Optimization with Application to
Communication and Queueing Systems}. Morgan \& Claypool, 2010.
\bibitem{nohadani18} O. Nohadani and K. Sharma. Optimization under decision-dependent uncertainty.
\emph{SIAM Journal on Optimization} 28(2):1773--1795, 2018.
\bibitem{nemhauser} G. L. Nemhauser, L. A. Wolsey, and M. L. Fisher. An analysis of approximations for
maximizing submodular set functions---I. \emph{Mathematical Programming} 14(1):265--294, 1978.
\bibitem{neuts67} M. F. Neuts. A general class of bulk queues with Poisson input. \emph{Annals of
Mathematical Statistics} 38(3):759--770, 1967.
\bibitem{ohtera90} H. Ohtera and S. Yamada. Optimal allocation and control problems for
software-testing resources. \emph{IEEE Transactions on Reliability} 39(2):171--176, 1990.
\bibitem{okati21} N. Okati, A. De, and M. Gomez-Rodriguez. Differentiable learning under triage.
2021. \url{https://arxiv.org/abs/2103.08902}.
\bibitem{onaga71} K. Onaga and O. Kakusho. On feasibility conditions of multicommodity flows in
networks. \emph{IEEE Transactions on Circuit Theory} 18(4):425--429, 1971.
\bibitem{ong24} I. Ong, A. Almahairi, V. Wu, W.-L. Chiang, T. Wu, J. E. Gonzalez, M. W. Kadous, and
I. Stoica. RouteLLM: learning to route LLMs with preference data. 2024.
\url{https://arxiv.org/abs/2406.18665}.
\bibitem{orlin18} J. B. Orlin, A. S. Schulz, and R. Udwani. Robust monotone submodular function
maximization. \emph{Mathematical Programming} 172(1--2):505--537, 2018.
\bibitem{pritsker69} A. A. B. Pritsker, L. J. Watters, and P. M. Wolfe. Multiproject scheduling with
limited resources: a zero-one programming approach. \emph{Management Science} 16(1):93--108, 1969.
\bibitem{raz86} T. Raz. A survey of models for allocating inspection effort in multistage production
systems. \emph{Journal of Quality Technology} 18(4):239--247, 1986.
\bibitem{soyster73} A. L. Soyster. Convex programming with set-inclusive constraints and applications
to inexact linear programming. \emph{Operations Research} 21(5):1154--1157, 1973.
\bibitem{tassiulas} L. Tassiulas and A. Ephremides. Stability properties of constrained queueing
systems and scheduling policies for maximum throughput in multihop radio networks. \emph{IEEE
Transactions on Automatic Control} 37(12):1936--1948, 1992.
\bibitem{williams63} A. C. Williams. Marginal values in linear programming. \emph{Journal of the
Society for Industrial and Applied Mathematics} 11(1):82--94, 1963.
\end{thebibliography}
\end{document}
